% CLOCKESCAPE10 | LSPV-3 | 2026-10-16
\documentclass{article}
\usepackage{amsmath,amsthm}
\newtheorem{theorem}{Theorem}
\begin{document}
\section{IDEAS PROPOSED BEFORE LITERATURE}
CANDIDATE A: recursion-theoretic fixed-point program might outrun a clocked self-output compiler; falsify with a program terminating outside the budget. CANDIDATE B: succinct status circuits might save encoding costs; test whether they also certify clock adequacy. CANDIDATE C: double-complement SAT compilers might yield an actual closed-formula fixed point; test distinct formula identities. CANDIDATE D: individual witness-program failures might yield one common SAT input; test covering portfolios. 
\section{LITERATURE CHECK AND REVISIONS}
Kleene (1938) supplies extensional program fixed points, not bounded-time equality. Cook (1971) reifies bounded computations with charged tableau size. Hartmanis--Stearns (1965) supplies time-hierarchy context; it does not imply SAT hardness. A yields the scoped theorem, B lacks clock adequacy, C fails on distinct sources, D lacks simultaneous-input quantifiers.
\section{CONSTRUCTION OR PROOF ATTEMPT}
Let $T(e)\ge1$ be total computable, and $E_T(e)$ an exact computable CNF whose satisfiability is equivalent to program $e$ emitting 1 within $T(e)$ machine steps. Ordinary existential tableau variables are not semantic global status bits.
\begin{theorem}[Clock escape]
There is a total constant-output fixed-point program $e_*$, with output $b$, such that $b=1-\operatorname{SAT}(E_T(e_*))$. Necessarily $\operatorname{SAT}(E_T(e_*))=0$, $b=1$, and the output occurs after $T(e_*)$.
\end{theorem}
\begin{proof}
The effective transformation sends $e$ to a program computing $1-\operatorname{SAT}(E_T(e))$, using an exact, possibly exponential SAT algorithm. Kleene's theorem gives $e_*$ with the same total constant-output behavior. If its bounded status were 1, it would emit 1 within the clock, but the fixed-point equality requires it to emit 0. Thus bounded status is 0, eventual output is 1 and the clock expires first.
\end{proof}
The compiler and tableau require polynomial work and length in $|e|+T(e)$; evaluating that tableau is not guaranteed inside $T(e)$, even under $P=NP$. This rules out a particular clock-adequacy inference, not polynomial SAT or P=NP. Distinct CNFs $F$, $J_D(F)$ and $J_D(J_D(F))$ alternate consistent SAT-status values. The lemma does not provide a satisfiable formula defeating all short witness programs on the same input.
\section{TESTS AND OUTCOME}
A reproducible local script checked 100 clocks $T=1,\ldots,100$, 303 assertions and zero failures; one-bit Boolean-map fixed-point histogram was (1,2,1). Test status COMPUTATIONAL. Full local manuscript and raw data are preserved separately. The smallest remaining target is a sound source-grounded realization theorem with charge-complete clock adequacy; P versus NP remains OPEN. No schedules modified.
\par\noindent\texttt{FINAL-SENTINEL-CLOCKESCAPE10-OPEN}
\end{document}